% ============================================================
%  MT-OPMNet — Complete LaTeX Figure Snippets (32 figures)
%  Copy the figures you need into your paper.
%
%  Prerequisites:
%    1. Run: python scripts/export_figures.py --format pdf
%    2. Copy figures/pdf/*.pdf into your LaTeX project
%    3. Add \usepackage{graphicx} to your preamble
% ============================================================


% ************************************************************
%  SECTION III — SYSTEM MODEL
% ************************************************************

% --- Fig 1: System Overview ---
\begin{figure*}[!t]
    \centering
    \includegraphics[width=\textwidth]{figures/pdf/system_overview.pdf}
    \caption{End-to-end architecture of the proposed MT-OPMNet optical performance monitoring system, from the WDM transmitter through the amplified fibre link to the coherent receiver and the multi-task inference engine.}
    \label{fig:system_model}
\end{figure*}

% --- Fig 2: Signal Processing Pipeline ---
\begin{figure}[!t]
    \centering
    \includegraphics[width=\columnwidth]{figures/pdf/signal_processing.pdf}
    \caption{Signal processing pipeline from optical transmission to MT-OPMNet monitoring outputs.}
    \label{fig:system_level}
\end{figure}

% --- Fig 3: AAH Examples for 5 Formats ---
\begin{figure*}[!t]
    \centering
    \includegraphics[width=\textwidth]{figures/pdf/results_aah_examples.pdf}
    \caption{Representative normalised AAHs for five modulation formats at OSNR\,=\,20\,dB and 28\,GBd.}
    \label{fig:aah_examples}
\end{figure*}

% --- Fig 4: AAH vs OSNR (16-QAM) ---
\begin{figure*}[!t]
    \centering
    \includegraphics[width=\textwidth]{figures/pdf/results_aah_osnr.pdf}
    \caption{Evolution of the 16-QAM AAH profile under varying OSNR conditions at 28\,GBd.}
    \label{fig:aah_osnr}
\end{figure*}

% --- Fig: Constellation Diagrams ---
\begin{figure*}[!t]
    \centering
    \includegraphics[width=\textwidth]{figures/pdf/constellations.pdf}
    \caption{Constellation diagrams of the five supported modulation formats: QPSK, 8-QAM, 16-QAM, 32-QAM, and 64-QAM.}
    \label{fig:constellations}
\end{figure*}


% ************************************************************
%  SECTION IV — PROPOSED FRAMEWORK
% ************************************************************

% --- Fig 5: MT-OPMNet Architecture ---
\begin{figure*}[!t]
    \centering
    \includegraphics[width=\textwidth]{figures/pdf/architecture.pdf}
    \caption{Detailed architecture of the proposed MT-OPMNet, comprising a shared four-block convolutional backbone, the channel-aware attention module (CAAM), and two task-specific output heads for OSNR regression and MFI classification.}
    \label{fig:architecture}
\end{figure*}

% --- Fig: CAAM Module ---
\begin{figure}[!t]
    \centering
    \includegraphics[width=\columnwidth]{figures/pdf/caam_module.pdf}
    \caption{Channel-Aware Attention Module (CAAM). Squeeze-and-excitation mechanism that learns channel-wise importance weights for adaptive feature recalibration.}
    \label{fig:caam}
\end{figure}

% --- Fig: Multi-Task Loss ---
\begin{figure}[!t]
    \centering
    \includegraphics[width=\columnwidth]{figures/pdf/multi_task_loss.pdf}
    \caption{Multi-task loss with homoscedastic uncertainty weighting. Smooth-L1 loss for OSNR regression and focal loss ($\gamma_f=2$) for MFI classification are combined using learnable log-variance parameters for automatic task balancing.}
    \label{fig:loss}
\end{figure}

% --- Fig: Training Pipeline ---
\begin{figure*}[!t]
    \centering
    \includegraphics[width=\textwidth]{figures/pdf/training_pipeline.pdf}
    \caption{Training and evaluation workflow of the MT-OPMNet framework, from dataset generation through model training to final testing.}
    \label{fig:training_workflow}
\end{figure*}


% ************************************************************
%  SECTION V — SIMULATION SETUP
% ************************************************************

% --- Fig 6: Dataset Distribution ---
\begin{figure}[!t]
    \centering
    \includegraphics[width=\columnwidth]{figures/pdf/results_dataset_distribution.pdf}
    \caption{Dataset sample distribution across the modulation format and OSNR range dimensions. Each cell contains 660 samples, confirming uniform coverage and the absence of class or range bias.}
    \label{fig:dataset_coverage}
\end{figure}


% ************************************************************
%  SECTION VI — RESULTS AND DISCUSSION
% ************************************************************

% --- Fig 7: OSNR Scatter Plot ---
\begin{figure}[!t]
    \centering
    \includegraphics[width=\columnwidth]{figures/pdf/results_osnr_scatter.pdf}
    \caption{Scatter plot of estimated versus true OSNR for the proposed MT-OPMNet. The tight clustering around the diagonal indicates accurate and unbiased predictions across the full OSNR range.}
    \label{fig:osnr_scatter}
\end{figure}

% --- Fig 8: OSNR Residual vs Distance ---
\begin{figure}[!t]
    \centering
    \includegraphics[width=\columnwidth]{figures/pdf/results_residual_distance.pdf}
    \caption{OSNR estimation residual ($\hat{y}_{\mathrm{OSNR}} - y_{\mathrm{OSNR}}$) as a function of transmission distance at 28\,GBd. The shaded band shows $\pm 1\sigma$.}
    \label{fig:residual_distance}
\end{figure}

% --- Fig 9: Confusion Matrix ---
\begin{figure}[!t]
    \centering
    \includegraphics[width=\columnwidth]{figures/pdf/results_confusion_matrix.pdf}
    \caption{Normalised confusion matrix for modulation format identification at 28\,GBd. Misclassifications are concentrated in a small number of off-diagonal cases.}
    \label{fig:confusion}
\end{figure}

% --- Fig: Top Misclassifications ---
\begin{figure}[!t]
    \centering
    \includegraphics[width=\columnwidth]{figures/pdf/results_top_confusions.pdf}
    \caption{Misclassification rates for the most frequently confused modulation format pairs at 28\,GBd.}
    \label{fig:top_confusions}
\end{figure}

% --- Fig 12: Per-Format OSNR RMSE ---
\begin{figure}[!t]
    \centering
    \includegraphics[width=\columnwidth]{figures/pdf/results_performat_rmse.pdf}
    \caption{Per-format OSNR estimation RMSE (dB) at 28\,GBd and 64\,GBd for the proposed MT-OPMNet. QPSK achieves the lowest error, while 64-QAM shows the highest RMSE.}
    \label{fig:performat}
\end{figure}

% --- Fig 13: MFI Accuracy vs OSNR ---
\begin{figure}[!t]
    \centering
    \includegraphics[width=\columnwidth]{figures/pdf/results_mfi_vs_osnr.pdf}
    \caption{MFI accuracy as a function of OSNR for four modulation formats at 28\,GBd. Higher-order formats are more sensitive to low OSNR conditions.}
    \label{fig:mfi_osnr}
\end{figure}

% --- Fig 14: RMSE vs Distance ---
\begin{figure}[!t]
    \centering
    \includegraphics[width=\columnwidth]{figures/pdf/results_rmse_vs_distance.pdf}
    \caption{OSNR estimation RMSE as a function of transmission distance at 28\,GBd for the evaluated models.}
    \label{fig:rmse_distance}
\end{figure}

% --- Fig 15: Ablation Study ---
\begin{figure}[!t]
    \centering
    \includegraphics[width=\columnwidth]{figures/pdf/results_ablation.pdf}
    \caption{Ablation study results: (a)~OSNR RMSE and (b)~MFI accuracy for four model configurations at 28\,GBd.}
    \label{fig:ablation}
\end{figure}

% --- Fig: Ablation Waterfall ---
\begin{figure}[!t]
    \centering
    \includegraphics[width=\columnwidth]{figures/pdf/results_ablation_waterfall.pdf}
    \caption{Ablation waterfall: cumulative OSNR RMSE improvement at 28\,GBd as architectural components are progressively added.}
    \label{fig:ablation_waterfall}
\end{figure}

% --- Fig 17: Radar Model Comparison ---
\begin{figure}[!t]
    \centering
    \includegraphics[width=\columnwidth]{figures/pdf/results_radar_comparison.pdf}
    \caption{Normalised multi-metric comparison of evaluated models. For each metric, higher scores indicate better performance. MT-OPMNet achieves the most balanced trade-off.}
    \label{fig:model_comparison}
\end{figure}

% --- Fig 18: OSNR Heatmap (Format x OSNR) ---
\begin{figure}[!t]
    \centering
    \includegraphics[width=\columnwidth]{figures/pdf/results_osnr_heatmap.pdf}
    \caption{Heatmap of OSNR estimation RMSE (dB) for each modulation format across five OSNR ranges at 28\,GBd. The mid-range OSNR region yields the lowest estimation errors.}
    \label{fig:osnr_heatmap}
\end{figure}

% --- Fig 19: Distance x Power Heatmap ---
\begin{figure}[!t]
    \centering
    \includegraphics[width=\columnwidth]{figures/pdf/results_dist_power_heatmap.pdf}
    \caption{OSNR RMSE heatmap across the transmission distance and per-channel launch power parameter space at 28\,GBd (16-QAM).}
    \label{fig:dist_power_heatmap}
\end{figure}

% --- Fig 20: Attention Weights + Bin Saliency ---
\begin{figure}[!t]
    \centering
    \includegraphics[width=\columnwidth]{figures/pdf/results_attention_saliency.pdf}
    \caption{Internal model analysis: (a)~learned CAAM attention weights across 256 feature channels for three modulation formats, and (b)~gradient-based saliency map showing bin importance for the OSNR regression task.}
    \label{fig:attention}
\end{figure}

% --- Fig: Task Weight Evolution ---
\begin{figure}[!t]
    \centering
    \includegraphics[width=\columnwidth]{figures/pdf/results_task_weights.pdf}
    \caption{Evolution of normalised task weights during training, showing adaptive shift from equal weighting towards OSNR-focused emphasis as MFI accuracy saturates.}
    \label{fig:task_weights}
\end{figure}

% --- Fig 22: Latency vs Batch Size ---
\begin{figure}[!t]
    \centering
    \includegraphics[width=\columnwidth]{figures/pdf/results_latency.pdf}
    \caption{Inference latency per sample versus batch size on the CPU platform. MT-OPMNet achieves consistently lower latency than the combined single-task baseline.}
    \label{fig:latency}
\end{figure}

% --- Fig 23: CDF of OSNR Errors ---
\begin{figure}[!t]
    \centering
    \includegraphics[width=\columnwidth]{figures/pdf/results_cdf.pdf}
    \caption{Cumulative distribution function (CDF) of the absolute OSNR estimation error at 28\,GBd. MT-OPMNet achieves the tightest error distribution among the compared models.}
    \label{fig:cdf}
\end{figure}

% --- Fig 24: Cross-Symbol-Rate Study ---
\begin{figure}[!t]
    \centering
    \includegraphics[width=\columnwidth]{figures/pdf/results_cross_rate.pdf}
    \caption{Cross-symbol-rate generalisation study: (a)~OSNR RMSE and (b)~MFI accuracy for four train/test symbol-rate configurations.}
    \label{fig:cross_rate}
\end{figure}

% --- Fig: Parameters + MACs ---
\begin{figure}[!t]
    \centering
    \includegraphics[width=\columnwidth]{figures/pdf/results_params_breakdown.pdf}
    \caption{Computational profile: (a)~MT-OPMNet parameter distribution across architectural components, and (b)~multiply-accumulate operations (MACs) comparison across all evaluated models.}
    \label{fig:params_macs}
\end{figure}

% --- Fig: OSNR Error Distribution ---
\begin{figure}[!t]
    \centering
    \includegraphics[width=\columnwidth]{figures/pdf/results_error_distribution.pdf}
    \caption{OSNR prediction error distribution. The near-zero mean confirms unbiased estimation.}
    \label{fig:error_dist}
\end{figure}

% --- Fig: OSNR Error per Modulation (Boxplot) ---
\begin{figure}[!t]
    \centering
    \includegraphics[width=\columnwidth]{figures/pdf/results_osnr_per_modulation.pdf}
    \caption{OSNR absolute error per modulation format. Higher-order modulations show slightly larger errors due to denser constellation spacing.}
    \label{fig:osnr_per_mod}
\end{figure}

% --- Fig: OSNR MAE vs True OSNR ---
\begin{figure}[!t]
    \centering
    \includegraphics[width=\columnwidth]{figures/pdf/results_osnr_vs_error.pdf}
    \caption{OSNR estimation error vs.\ OSNR level. Performance improves at higher OSNR values where signal-to-noise separation is clearer.}
    \label{fig:osnr_vs_error}
\end{figure}

% --- Fig: Training Convergence (full-width) ---
\begin{figure*}[!t]
    \centering
    \includegraphics[width=\textwidth]{figures/pdf/results_training_curves.pdf}
    \caption{Training convergence curves: (a)~total loss, (b)~MFI classification accuracy, and (c)~validation OSNR MAE over training epochs.}
    \label{fig:convergence}
\end{figure*}
